\section{Training-Stage Results}
\label{sec:training}

\subsection{Block A: Pretraining}
\label{sec:training:pretrain}

Table~\ref{tab:pretrain_dense} reports pretraining throughput and MFU for the
dense configuration across GPU counts and vendors.
Table~\ref{tab:pretrain_moe} gives the corresponding MoE results.

\begin{table}[h]
\centering
\caption{Dense pretrain (hidden\_size=768, 8 layers, bf16, 2 epochs, pretrain\_t2t\_mini).
MFU computed against H100 989\,TFLOPs/s (NVIDIA) and MI350 \PH{X}\,TFLOPs/s (AMD).}
\label{tab:pretrain_dense}
\small
\begin{tabular}{lllrrrrr}
\toprule
\textbf{Vendor} & \textbf{GPUs} & \textbf{dtype / compile} &
\textbf{tok/s} & \textbf{MFU (\%)} & \textbf{Peak HBM (GB)} &
\textbf{Epoch (min)} & \textbf{Final loss} \\
\midrule
H100  & 1 & bf16 / off & \PH{} & \PH{} & \PH{} & \PH{} & \PH{} \\
H100  & 4 & bf16 / off & \PH{} & \PH{} & \PH{} & \PH{} & \PH{} \\
H100  & 8 & bf16 / off & \PH{} & \PH{} & \PH{} & \PH{} & \PH{} \\
H100  & 8 & bf16 / on  & \PH{} & \PH{} & \PH{} & \PH{} & \PH{} \\
H100  & 8 & fp16 / off & \PH{} & \PH{} & \PH{} & \PH{} & \PH{} \\
\midrule
MI350 & 1 & bf16 / off & \PH{} & \PH{} & \PH{} & \PH{} & \PH{} \\
MI350 & 4 & bf16 / off & \PH{} & \PH{} & \PH{} & \PH{} & \PH{} \\
MI350 & 8 & bf16 / off & \PH{} & \PH{} & \PH{} & \PH{} & \PH{} \\
MI350 & 8 & bf16 / on  & \PH{} & \PH{} & \PH{} & \PH{} & \PH{} \\
MI350 & 8 & fp16 / off & \PH{} & \PH{} & \PH{} & \PH{} & \PH{} \\
\bottomrule
\end{tabular}
\end{table}

\begin{table}[h]
\centering
\caption{MoE pretrain (4 experts, 1 active, hidden\_size=768, bf16). MFU uses $P_{\text{active}}$.}
\label{tab:pretrain_moe}
\small
\begin{tabular}{llrrrrr}
\toprule
\textbf{Vendor} & \textbf{GPUs} &
\textbf{tok/s} & \textbf{MFU (\%)} & \textbf{Peak HBM (GB)} &
\textbf{Epoch (min)} & \textbf{Final loss + aux} \\
\midrule
H100  & 1 & \PH{} & \PH{} & \PH{} & \PH{} & \PH{} \\
H100  & 4 & \PH{} & \PH{} & \PH{} & \PH{} & \PH{} \\
H100  & 8 & \PH{} & \PH{} & \PH{} & \PH{} & \PH{} \\
\midrule
MI350 & 1 & \PH{} & \PH{} & \PH{} & \PH{} & \PH{} \\
MI350 & 4 & \PH{} & \PH{} & \PH{} & \PH{} & \PH{} \\
MI350 & 8 & \PH{} & \PH{} & \PH{} & \PH{} & \PH{} \\
\bottomrule
\end{tabular}
\end{table}

\paragraph{Multi-GPU scaling.}
Figure~\ref{fig:scaling} (placeholder) shows strong-scaling efficiency
(normalized throughput per GPU vs.\ ideal linear) for both platforms.
We expect near-linear scaling within 8 GPUs for a 64\,M-parameter model
whose activation size is small relative to intra-node interconnect bandwidth.

\paragraph{\texttt{torch.compile} effect.}
\PH{Discuss compile speedup on H100 vs MI350 once data are available.}

\subsection{Block B: Supervised Fine-Tuning, LoRA, and Distillation}
\label{sec:training:posttrain}

\begin{table}[h]
\centering
\caption{Post-training stages: throughput and memory. SFT uses seq\_len=768,
LoRA trains only adapter parameters ($\ll 1\%$ of total), Distillation
runs MoE teacher + dense student simultaneously.}
\label{tab:posttrain}
\small
\begin{tabular}{lllrrr}
\toprule
\textbf{Stage} & \textbf{Vendor} & \textbf{GPUs} &
\textbf{tok/s} & \textbf{Peak HBM (GB)} & \textbf{Final loss} \\
\midrule
SFT         & H100  & 1 & \PH{} & \PH{} & \PH{} \\
SFT         & H100  & 8 & \PH{} & \PH{} & \PH{} \\
SFT         & MI350 & 1 & \PH{} & \PH{} & \PH{} \\
SFT         & MI350 & 8 & \PH{} & \PH{} & \PH{} \\
\midrule
LoRA        & H100  & 1 & \PH{} & \PH{} & \PH{} \\
LoRA        & H100  & 8 & \PH{} & \PH{} & \PH{} \\
LoRA        & MI350 & 1 & \PH{} & \PH{} & \PH{} \\
LoRA        & MI350 & 8 & \PH{} & \PH{} & \PH{} \\
\midrule
Distill (MoE$\to$dense) & H100  & 8 & \PH{} & \PH{} & \PH{} \\
Distill (MoE$\to$dense) & MI350 & 8 & \PH{} & \PH{} & \PH{} \\
\bottomrule
\end{tabular}
\end{table}

\paragraph{LoRA efficiency.}
Because \texttt{torch.compile} is automatically disabled for LoRA (the
monkey-patched forward is incompatible; see \texttt{train\_lora.py:165}),
the throughput gap vs.\ full SFT is entirely a function of trainable-parameter
ratio, not kernel compilation.

\paragraph{Distillation with a MoE teacher.}
The distillation step exercises MoE forward passes in the teacher alongside
dense backward passes in the student on every step, stressing memory bandwidth
differently from pure-MoE or pure-dense training.
\PH{Discuss observed memory and throughput differences between vendors.}

\subsection{Block C: Preference Alignment (DPO, PPO, GRPO)}
\label{sec:training:rlhf}

\begin{table}[h]
\centering
\caption{DPO results. The frozen reference doubles activation memory;
 batch size is limited to 4 per GPU.}
\label{tab:dpo}
\small
\begin{tabular}{llrrrr}
\toprule
\textbf{Vendor} & \textbf{GPUs} &
\textbf{tok/s} & \textbf{Peak HBM (GB)} &
\textbf{DPO loss} & \textbf{Epoch (min)} \\
\midrule
H100  & 1 & \PH{} & \PH{} & \PH{} & \PH{} \\
H100  & 8 & \PH{} & \PH{} & \PH{} & \PH{} \\
MI350 & 1 & \PH{} & \PH{} & \PH{} & \PH{} \\
MI350 & 8 & \PH{} & \PH{} & \PH{} & \PH{} \\
\bottomrule
\end{tabular}
\end{table}

\begin{table}[h]
\centering
\caption{GRPO and PPO rollout throughput and training metrics.
Rollout-tok/s is measured inside the training loop via the \texttt{RolloutResult} timer.}
\label{tab:rlhf}
\small
\begin{tabular}{lllrrrrr}
\toprule
\textbf{Algo} & \textbf{Vendor} & \textbf{Engine} &
\textbf{$G$} & \textbf{rollout-tok/s} & \textbf{train-tok/s} &
\textbf{Peak HBM (GB)} & \textbf{Reward} \\
\midrule
GRPO & H100  & torch  & 4 & \PH{} & \PH{} & \PH{} & \PH{} \\
GRPO & H100  & torch  & 8 & \PH{} & \PH{} & \PH{} & \PH{} \\
GRPO & H100  & sglang & 4 & \PH{} & \PH{} & \PH{} & \PH{} \\
GRPO & MI350 & torch  & 4 & \PH{} & \PH{} & \PH{} & \PH{} \\
GRPO & MI350 & torch  & 8 & \PH{} & \PH{} & \PH{} & \PH{} \\
GRPO & MI350 & sglang & 4 & \PH{} & \PH{} & \PH{} & \PH{} \\
\midrule
PPO  & H100  & torch  & -- & \PH{} & \PH{} & \PH{} & \PH{} \\
PPO  & H100  & sglang & -- & \PH{} & \PH{} & \PH{} & \PH{} \\
PPO  & MI350 & torch  & -- & \PH{} & \PH{} & \PH{} & \PH{} \\
PPO  & MI350 & sglang & -- & \PH{} & \PH{} & \PH{} & \PH{} \\
\bottomrule
\end{tabular}
\end{table}

\subsection{Block D: Agent RL}
\label{sec:training:agent}

Table~\ref{tab:agent} reports trajectory throughput and tool-call success rate.
Because agent training involves multi-turn rollouts with real tool invocations,
the rollout step is significantly more variable than in GRPO/PPO.

\begin{table}[h]
\centering
\caption{Agent RL: trajectory throughput and tool-call success.}
\label{tab:agent}
\small
\begin{tabular}{llrrrr}
\toprule
\textbf{Vendor} & \textbf{GPUs} &
\textbf{traj/h} & \textbf{tool-call success (\%)} &
\textbf{Peak HBM (GB)} & \textbf{Reward curve} \\
\midrule
H100  & 1 & \PH{} & \PH{} & \PH{} & \PH{} \\
H100  & 8 & \PH{} & \PH{} & \PH{} & \PH{} \\
MI350 & 1 & \PH{} & \PH{} & \PH{} & \PH{} \\
MI350 & 8 & \PH{} & \PH{} & \PH{} & \PH{} \\
\bottomrule
\end{tabular}
\end{table}
